\abstract

Ключевые слова: распределённые системы, межкластерная репликация,
геораспределённая архитектура, отказоустойчивость,
репликационный журнал, кворум, стохастическая модель,
параметрическая оптимизация.

Данная работа посвящена исследованию геораспределённой системы
межкластерной репликации данных в архитектуре «активный–пассивный».
Рассматривается модель распределённого хранилища,
состоящего из двух кластеров,
внутри которых согласование операций осуществляется
на основе кворумного механизма,
а между кластерами реализуется асинхронная передача
репликационного журнала.

Целью работы является формализация задачи обеспечения
целостности, доступности и согласованности данных
в условиях случайных сетевых задержек,
потерь сообщений и отказов узлов.
Для этого вводится вероятностная модель функционирования системы,
определяются количественные критерии качества,
а также формулируется задача выбора параметров
отказоустойчивости и репликации
в виде задачи оптимизации при вероятностных ограничениях.

Предлагаемый подход позволяет определить допустимую область
настроек системы, обеспечивающих выполнение требований
по целостности, доступности и согласованности данных,
и выбрать конфигурацию параметров,
минимизирующую эксплуатационные издержки.
Полученные результаты могут быть использованы
при проектировании и настройке
геораспределённых систем хранения данных.
